\chapter{Introdução}

A proteção de um ambiente computacional combina controles em camadas distintas: a
forma como a rede é endereçada e segmentada, as regras que filtram o tráfego entre os
hosts e, nos serviços de nuvem, quem pode executar quais ações sobre quais recursos. Uma
falha em qualquer dessas camadas costuma bastar para expor dados, como mostram os
recorrentes casos de \textit{buckets} de armazenamento deixados públicos por engano.

Este relatório documenta um laboratório controlado que exercita essas três camadas. Duas
máquinas virtuais Ubuntu Server 22.04, ligadas por uma rede interna isolada do VirtualBox,
fazem os papéis de \emph{servidor} e \emph{cliente}. Sobre elas foram realizadas quatro
fases: (1) endereçamento estático e divisão da rede 192.168.50.0/24 em sub-redes /26;
(2) publicação de um serviço web na porta 8080 e de SSH, com um \textit{firewall} local
(\texttt{ufw}) que libera somente a porta 8080 vinda do cliente, validado com
\texttt{nmap} antes e depois; (3) criação de um usuário com permissão mínima sobre um
\textit{bucket} compatível com S3 e a correção proposital desse \textit{bucket} de público
para privado; e (4) uma verificação final que reúne evidências de todos os controles.

Como o projeto MinIO de código aberto, sugerido pelo enunciado para simulação local, foi
arquivado e teve os binários retirados do site oficial, a Fase 3 utilizou o SILO, um
\textit{fork} comunitário mantido do MinIO que preserva a mesma API S3, a mesma API
administrativa e o mesmo cliente de linha de comando (renomeado \texttt{mcli})
\cite{silo-docs,minio-arquivado}.

Todas as etapas foram automatizadas por \textit{scripts} que registram cada comando,
sua saída e seu código de saída; esses registros compõem as evidências apresentadas no
\autoref{cap:desenvolvimento} e no Anexo~\ref{anexo:roteiro}. O documento segue a
NBR 14724 \cite{abnt14724}.
